% Section 3. Sources: paper2/problem_transformations.md (problems, statements,
% proofs), paper2/equivalences.md (master table, sources, gaps),
% paper2/proof_reductions.md (faults in published proofs). Lean:
% lean/MOSPFormalization/ and Complex/; axioms printed by paper2/axiom_check.lean.
\section{The equivalences, proved}
\label{sec:equiv}

\citet[p.~1764]{LY2002} introduce Table~1 as ``a set of problems that consist
of, given input $\Pi$, compute a function $f(\Pi)$ that is either equal to the
number of open stacks or closely related to it (plus or minus one)''. The table
cites where each problem was studied. It gives no proofs and no graph. Two
readings of ``plus or minus one'' must be kept apart. Under the \emph{literal}
reading, $f$ lies within one of $Z$, the number of open stacks of the
corresponding instance, on every input. Under the \emph{fixed-offset} reading,
$f = Z + c$ for one constant $c\in\{-1,0,1\}$ on every input. Only the second
makes computing one problem the same as computing another.

This section states each problem as an instance and a question, gives the
relation to pathwidth that is true, and says where it is proved. Every relation
stated below is a theorem of our Lean development unless it says otherwise. All
the theorems named depend only on the axioms \texttt{propext},
\texttt{Classical.choice} and \texttt{Quot.sound}.\footnote{Printed by
\texttt{lake env lean paper2/axiom\_check.lean}, which covers every theorem
named in this paper. The development has one
\texttt{sorry}, a conjecture about lower bounds that no result of this paper
uses.} Table~\ref{tab:master} summarises the result and
Figure~\ref{fig:chain} draws it.

\subsection{The problems}
\label{sec:problems}

Each problem below is a decision problem with an instance and an integer $k$;
its value is the least $k$ for which the answer is yes.

\begin{description}
\item[Pathwidth $\pw(G)$.] Is there a sequence of bags $X_1,\dots,X_r\subseteq V$
  covering $V$, with every edge inside some bag, $X_i\cap X_\ell\subseteq X_j$
  whenever $i\le j\le\ell$, and $|X_i|\le k+1$? \citep[p.~346]{LY2002ref13}.
\item[Vertex separation $\vs(G)$.] Is there a layout $L$ such that for every
  $i<n$ at most $k$ vertices $u$ with $L(u)\le i$ have a neighbour $v$ with
  $L(v)>i$? \citep[p.~346]{LY2002ref13}.
\item[MOSP $Z(M)$.] Given a 0/1 matrix $M$ of customers by patterns, is there
  an order of the patterns with at most $k$ customers active at every position?
  \citep[eqs.~(1)--(2)]{LY2002}.
\item[Gate matrix layout $t(M)$.] Given a net--gate matrix, is there an order of
  the gates and an assignment of nets to $k$ tracks such that nets active at a
  common position get different tracks? \citep[p.~18]{LY2002ref6};
  \citep[Problem~1]{LY2002ref8}.
\item[One-dimensional logic.] Given gates, nets and the nets each gate connects,
  is there an order of the gates whose net intervals have an interval graph of
  clique number at most $k$? Its \emph{connection graph} $H$ joins two nets that
  share a gate. A \emph{boundary} variant pins two gates to the ends
  \citep[\S\S II--IV]{LY2002ref7}.
\item[PLA folding $\pla(M)$.] As gate matrix layout, with at most two nets per
  track. With no such bound (\emph{multiple folding}) it is gate matrix layout
  itself \citep[pp.~25, 36]{LY2002ref6}.
\item[Interval thickness $\theta(G)$.] Is there an interval supergraph of $G$ on
  the same vertices with clique number at most $k$?
  \citep[p.~182]{LY2002ref9}; \citep[p.~31]{LY2002ref6}.
\item[Node search $\ns(G)$.] Can searchers, placed on and removed from vertices,
  clear every edge with at most $k$ on the graph at once, where an edge is
  cleared when both ends are guarded and a clear edge is recontaminated by a
  searcher-free path to a contaminated one? $\mns(G)$ is the same over strategies
  that never recontaminate \citep[p.~181]{LY2002ref9}; \citep[\S2]{LY2002ref10}.
\item[Edge search $\es(G)$.] The same game with a third move, sliding a searcher
  along an edge, which clears that edge; $\pes(G)$ for progressive strategies
  \citep[p.~208]{LY2002ref10}; \citep[p.~53]{EllisSudboroughTurner1994}.
\item[Narrowness $\nu(G)$.] Is there an order $v_1,\dots,v_n$ such that the
  following process never holds more than $k$ vertices: at step $i$ add $v_i$,
  then remove every $v_j$, $j\le i$, with no neighbour $v_\ell$, $\ell>i$?
  \citep[\S2]{LY2002ref11}.
\item[Split bandwidth $\sbw(G)$.] Is there a graph obtained from $G$ by
  repeatedly splitting a vertex into an edge, its neighbours partitioned between
  the two ends, with bandwidth at most $k$? \citep[\S3.2]{LY2002ref12}.
\item[Edge separation.] Table~1 cites \citet{LY2002ref14}, which supports three
  readings: cutwidth $\cw(G)$, the least over layouts of the most edges crossing
  a gap (p.~468); modified cutwidth $\mcw(G)$, counting only edges that pass
  strictly over a vertex (Definition~6); and Lengauer's vertex separator game
  $\VSG(G)$, the least $K$ such that the vertices can be pebbled in an order that
  leaves at most $K$ unpebbled vertices adjacent to pebbled ones (p.~467).
\end{description}

\subsection{What is true}

\begin{table}[t]
\centering\small
\caption{Table~1 of \citet{LY2002}, settled. \emph{Relation}: what is true and
proved in Lean, on the graph the problem lives on (the MOSP graph or net graph
for the matrix problems, the input graph otherwise). \emph{Status}:
\emph{exact}, a fixed offset; \emph{band}, an inequality whose ends both occur;
\emph{false}, Table~1's claim fails, with a counterexample family proved in
Lean.}
\label{tab:master}
\begin{tabularx}{\linewidth}{@{}>{\raggedright}p{2.7cm}>{\raggedright}p{3.4cm}>{\raggedright\arraybackslash}X l@{}}
\toprule
Problem & Proved in & Relation & Status \\
\midrule
MOSP & \citet{Yanasse1997a}; \citet{FellowsLangston1989}; \citet{LY2002} &
  $Z(M)=\pw(G_M)+1$ ($M$ has a 1) & exact \\
Gate matrix layout & \citet{LY2002ref6} & $t(M)=Z(M)=\pw(G_M)+1$ & exact \\
One-dim.\ logic & \citet{LY2002ref7} & tracks $=\theta(H)=\pw(H)+1$; boundary variant $\ge\pw(H)+1$ & exact \\
PLA folding & \citet{LY2002ref6} & $\max(\pw+1,\lceil|N|/2\rceil)\le\pla$, gap unbounded; multiple folding $=\pw+1$ & false; variant exact \\
Interval thickness & \citet{LY2002ref6}; \citet{LY2002ref9} & $\theta=\pw+1$ ($V\ne\emptyset$) & exact \\
Node search & \citet{LY2002ref9,LY2002ref10}; \citet{BienstockSeymour1991} & $\ns=\mns=\vs+1=\pw+1=\theta$ ($E\ne\emptyset$) & exact \\
Edge search & \citet{EllisSudboroughTurner1994}; \citet{LY2002ref10} & $\vs\le\es\le\vs+2$, also for $\pes$ & band \\
Narrowness & \citet{LY2002ref11} & $\nu=\pw+1$ ($V\ne\emptyset$) & exact \\
Split bandwidth & \citet{LY2002ref12} & $\pw\le\sbw\le\pw+1$ & band \\
Path-width & \citet{LY2002ref13} & the reference & definition \\
Edge separation & \citet{LY2002ref14} & $\cw,\mcw$ unbounded against $\pw$ (stars); $\VSG=\max(1,\vs)$ & false; VSG exact \\
Vertex separation & \citet{LY2002ref13} & $\vs=\pw$ & exact \\
\bottomrule
\end{tabularx}
\end{table}

\begin{figure}[t]
\centering
\includegraphics[width=\linewidth]{sec3_fig1_chain}
\caption{The equivalences of Table~\ref{tab:master} as a graph around
pathwidth. Solid edges are exact relations, dashed edges bands, both proved in
Lean; the dotted edge is LaPaugh's $\mathrm{es}=\mathrm{pes}$, the one relation
not proved, which no row needs. Dark red marks Table~1's two false claims, each
refuted by a counterexample family proved in Lean. The four matrix problems (top
left) live on the MOSP graph $G_M$ or the net graph $H$; $\mathrm{mpb}$ is the minimum
progressive pebbling number of Section~\ref{sec:thirteenth}. Hypotheses are
omitted: $M$ has a 1, $V\neq\emptyset$, or $G$ has an edge, as in
Table~\ref{tab:master}. Two edges that follow from the others,
$\mathrm{mns}=\theta$ and $\mathrm{ns}=\mathrm{vs}+1$, are not drawn.}
\label{fig:chain}
\end{figure}

Two lemmas carry most of the exact rows. For a path decomposition and a vertex
$v$, let $\mathrm{first}(v)$ and $\mathrm{last}(v)$ be the least and greatest
indices of bags containing $v$; the decomposition's interval property says
$v\in X_i$ exactly when $\mathrm{first}(v)\le i\le\mathrm{last}(v)$.

\begin{lemma}[Helly for intervals]\label{lem:helly}
Pairwise intersecting intervals of a line have a common point.
\end{lemma}
\begin{proof}
Let $a$ be the largest left end. Any two intervals meet, so every right end is
at least every left end, and in particular at least $a$. So $a$ lies in all of
them.
\end{proof}

\begin{lemma}[cliques sit in a bag]\label{lem:clique}
In a path decomposition of $G$, every clique of $G$ lies inside one bag.\lean{PathDecomposition.exists_bag_of_isClique}
\end{lemma}
\begin{proof}
The bags containing a vertex $u$ of the clique form the interval
$[\mathrm{first}(u),\mathrm{last}(u)]$. Two vertices of the clique are
adjacent, so their edge lies in a bag and their intervals meet. By
Lemma~\ref{lem:helly} all the intervals share an index $i$, and the clique lies
in $X_i$.
\end{proof}

\begin{theorem}[vertex separation is pathwidth]\label{thm:vs}
For every graph, $\vs(G)=\pw(G)$.\lean{vertexSeparation_eq_pathwidth}
\end{theorem}
This is Theorem~3.1 of \citet{LY2002ref13}. From a layout $L$ of separation $k$,
the bags $X_i = V_L(i-1)\cup\{L^{-1}(i)\}$, where $V_L(i)$ is the set the
definition counts, form a decomposition of width $k$. Conversely, ordering the
vertices by $\mathrm{first}$ gives a layout in which everything counted at a cut
lies in the bag where the next vertex first appears, beside that vertex.

\begin{theorem}[MOSP is pathwidth plus one]\label{thm:mosp}
If $M$ has at least one 1, then $Z(M)=\pw(G_M)+1$.\lean{MOSPInstance.mospValue_eq_pathwidth_add_one}
\end{theorem}
\begin{proof}
($\le$) Take a decomposition of $G_M$ of width $k$. The rows of each column form
a clique, so by Lemma~\ref{lem:clique} they lie in a bag $X_{b(c)}$. Order the
columns by $b(c)$. If row $r$ is active at the position of column $c$, it has
columns $x,y$ with $b(x)\le b(c)\le b(y)$, and $r\in X_{b(x)}\cap X_{b(y)}$, so the
interval property gives $r\in X_{b(c)}$. At most $k+1$ rows are active there.

($\ge$) Take an order with at most $Z$ active rows at every position. The
active sets, with a singleton bag for each row with no 1, form a path
decomposition: rows sharing a column are both active at its position, and a row
is active on an interval of positions. Its width is $Z-1$, using $Z\ge1$.
\end{proof}

Read as a cutting problem, Theorem~\ref{thm:mosp} is Proposition~5 of
\citet{Yanasse1997a}, which turns each pattern into a clique on its piece types;
read as a layout problem, it is Theorem~7 of \citet{FellowsLangston1989}, with
the column-expansion Lemma~4.1 of \citet{FellowsLangston1987}; and
\citet[Prop.~2]{LY2002} identify MOSP with gate matrix layout. The proof above
needs the Helly property and the interval property, not Hall's theorem. It needs
no assumption that the instance is reduced.

The graph matters. The pattern graph, with the patterns as vertices joined when
they share a customer, gives no bound in either direction: one pattern shared by
four customers who each have a private pattern has $Z=4$, while its pattern
graph is the star $K_{1,4}$, of pathwidth~1.\lean{star_mospValue, star_pathwidth_agreementGraph, star_refutes_pattern_graph_bound} A single
customer with many patterns errs the other way.

The other exact rows attach to Theorems~\ref{thm:vs} and~\ref{thm:mosp}.

\begin{theorem}[the exact core]\label{thm:core}
\leavevmode
\begin{enumerate}
\item Gate matrix layout is MOSP on the same matrix: if $M$ has a 1,
  $t(M)=Z(M)=\pw(G_M)+1$. For a fixed gate order the fewest tracks equals the
  most nets active at one position, by the left-edge algorithm.\lean{NetGateMatrix.tracks_eq_mospValue, NetGateMatrix.tracks_eq_pathwidth_add_one}
\item Interval thickness: if $V\ne\emptyset$, $\theta(G)=\pw(G)+1$
  \citep[Prop.~3.5]{LY2002ref6}.\lean{intervalThickness_eq_pathwidth_add_one}
\item One-dimensional logic: if every net lies on a gate, the fewest tracks is
  $\theta(H)$, and if some gate connects a net it is $\pw(H)+1$
  \citep[Thm~3]{LY2002ref7}.\lean{LogicArray.tracks_eq_intervalThickness, LogicArray.tracks_eq_pathwidth_add_one}
\item Narrowness: if $V\ne\emptyset$, $\nu(G)=\pw(G)+1$, and for each order
  $\sigma$, $\nu(\sigma)=\vs(\sigma^{\mathrm{rev}})+1$
  \citep[Prop.~3.1]{LY2002ref11}.\lean{narrowness_eq_pathwidth_add_one}
\item Node search: for every graph $\ns(G)=\mns(G)$, and if $E\ne\emptyset$,
  $\ns(G)=\vs(G)+1=\pw(G)+1=\theta(G)$
  \citep[Thm]{LY2002ref9}; \citep[Thms~2.3, 4.1]{LY2002ref10}.\lean{nodeSearchMonotonicity, nodeSearch_eq_vertexSeparation_add_one, nodeSearch_chain}
\item Multiple folding: the fewest tracks is $t(M)=\pw(G_M)+1$
  \citep[Thm~3.14]{LY2002ref6}.\lean{NetGateMatrix.foldTracks_eq_pathwidth_add_one}
\item Lengauer's vertex separator game: $\VSG(G)=\max(1,\vs(G))$
  \citep[Thm~4]{LY2002ref14}.\lean{vsg_eq_max}
\end{enumerate}
\end{theorem}

The Lean proof that node search is monotone follows the crusade argument of
\citet{BienstockSeymour1991}, not the reduction to edge search of
\citet{LaPaugh1993} on which \citet{LY2002ref10} rely.

\begin{theorem}[bands]\label{thm:bands}
For every graph, $\pw(G)\le\sbw(G)\le\pw(G)+1$ \citep[Thm~8]{LY2002ref12}, and
$\vs(G)\le\es(G)\le\vs(G)+2$, and the same for $\pes$
\citep[Thm~2.1]{EllisSudboroughTurner1994}.\lean{pathwidth_le_splitBandwidth_le_pathwidth_add_one, vertexSeparation_le_edgeSearch_le_add_two, vertexSeparation_le_progressiveEdgeSearch_le_add_two}
\end{theorem}
Every end of both bands occurs. $K_2$ has $\sbw=\pw=1$ and $\es=\vs=1$; the star
$K_{1,3}$ has $\pw=1$ and $\sbw=2$, and $\vs=1$ with $\es=2$; and $K_{3,3}$ has
$\vs=3$ with $\es=5$. These values come from a brute-force check of every
quantity from its own definition, not from Lean.\footnote{\texttt{python -m
paper2.complex\_check}; $\sbw(K_{1,3})\ge2$ is also argued by hand: splits of a
tree are trees and keep at least three leaves.} Fomin assumes a connected graph
with at least two vertices; the Lean proof needs neither. The equality
$\es=\pes$ \citep{LaPaugh1993} is stated in Lean as a named proposition and used
by no row.

\begin{theorem}[two rows are false]\label{thm:false}
\leavevmode
\begin{enumerate}
\item PLA folding: for a matrix with a 1,
  $\max(\pw(G_M)+1,\lceil|N|/2\rceil)\le\pla(M)$, and the gap to $\pw+1$ is
  unbounded. The $5\times5$ identity needs 3 tracks while $\pw(G_M)+1=1$, and the
  path $P_n$ as a matrix, one gate per edge, has $t=2$ and
  $\pla\ge\lceil n/2\rceil$, so the gap is unbounded on connected instances too
  \citep[Prop.~3.15]{LY2002ref6}.\lean{plaTracks_idMatrix_five, plaTracks_idMatrix_unbounded, plaTracks_pathMatrix_unbounded}
\item Edge separation read as cutwidth: $\pw\le\cw$ always, but
  $\cw(K_{1,n})\ge n/2$ and $\mcw(K_{1,n})\ge n/2-1$ while $\pw(K_{1,n})=1$; so
  neither cutwidth reading is within a constant of pathwidth.\lean{cutwidth_unbounded, modCutwidth_unbounded}
\end{enumerate}
\end{theorem}
Table~1's PLA claim holds for multiple folding (Theorem~\ref{thm:core}). The
edge-separation row is also misattributed: what \citet{LY2002ref14} proves
exactly is the vertex separator game, which is Table~1's vertex separation row.

\subsection{Edge cases and a gap in a published proof}

Three of the published equalities fail on degenerate inputs, and each failure
is proved in Lean. On a nonempty graph with no edges, $\ns=0$ while
$\vs+1=\theta=1$, so both equalities of \citet{LY2002ref9,LY2002ref10} need an
edge.\lean{nodeSearch_ne_intervalThickness_of_edgeless, monotoneNodeSearch_ne_vertexSeparation_add_one_of_edgeless}
For a matrix with no 1s, $Z=0$ while $\pw+1=1$. Lengauer's Theorem~4 fails at
$K=0$ on an edgeless graph.\lean{isPositiveVSG_triangleGraph_counterexample}
\citet[Thm~3.9]{LY2002ref6}, $t=\ns$, fails literally on the $1\times1$ matrix
$[1]$.\lean{tracks_ne_monotoneNodeSearch_one}

The published proof of $\vs\le\ns-1$ \citep[Thm~4.1]{LY2002ref10} has a gap. It
orders the vertices by when they first receive a searcher and claims that every
vertex counted at the worst cut still holds a searcher, ``since the vertices in
$D_{i_0}$ are joined to vertices not yet cleared'' (p.~217). A strategy may place
a searcher and remove it again before clearing anything, which recontaminates
nothing. On $K_{1,3}$, placing and removing a searcher on each leaf before
guarding the centre is an optimal recontamination-free strategy for which the
claim fails.\lean{kirousisPapadimitriou_claim2_false} The theorem is true. The
repair orders the vertices by clearing time, as our Lean proof does, or uses the
paper's own normal form, its Corollary~2.4.

\subsection{A thirteenth member}\label{sec:thirteenth}

Progressive black-white pebbling, which is not in Table~1, belongs to the exact
core. \citet[Thm~3]{LY2002ref14} gives
$\vs(G)\le K\iff\mathrm{pbw}(G_d)\le K+2$ on the dag $G_d$ built from $G$; his
Theorem~2 gives $\mathrm{pbw}(D)=\vs(D_u)+1$ on a dag $D$, where $D_u$ is the
MOSP graph of the matrix whose column $v$ is $v$ with its predecessors; and the
minimum progressive pebbling number of \citet[Thm~3.1]{LY2002ref10} equals
$\pw+1$. All three are proved in Lean, each with its hypothesis: Lengauer's
Theorem~2 in its instance form fails at $K=1$ on a dag with no arcs, and
$\mathrm{mpb}=\ns$ fails on edgeless graphs.\lean{pbw_lengauerD_eq_pathwidth, pbw_eq_mospValue_pebbleMatrix, mpb_eq_pathwidth_add_one}

\paragraph{Summary.} Seven of Table~1's twelve problems are pathwidth, or
pathwidth plus one, exactly: MOSP, gate matrix layout, one-dimensional logic,
interval thickness, node search, narrowness and vertex separation, with pathwidth
itself as the reference. So are two variants, multiple folding and Lengauer's
vertex separator game. Split bandwidth and edge search are bands of width one
and two, which meet the literal reading of Table~1 and not the fixed offset.
Simple PLA folding and edge separation are false as stated. One theorem is
stated and not proved, LaPaugh's $\es=\pes$, and no row needs it. One of
Table~1's references, Kashiwabara and Fujisawa (1979), is not held; its row is
proved from \citet{LY2002ref6} and \citet{LY2002ref9}.
